\section{Discussion}
\label{sec:discussion}

\subsection{Honest novelty, restated against the prior art}
\label{sec:novelty}

We fix the boundary once more, now that the results are on the table, because the value of
this paper is its discipline about what is and is not new.

\emph{Versus rigidity-based autonomous fault monitoring \citep{iiyama2024,iiyama2025edm}.} The
recent rigidity and Euclidean-distance-matrix methods are the closest prior art: they monitor
an inter-satellite-linked constellation autonomously, and the EDM follow-on already detects
clock-jump faults and identifies a \emph{single} faulty satellite. Three distinctions remain
genuinely ours. (a)~Theirs is a simulation-validated \emph{algorithm}---a rigidity/GCEDM
singular-value test---whereas ours is an \emph{algebraic characterization} of which faults are
provably invisible (undetectable $\iff\rng(\Gmat)$) with provable bounds. (b)~They identify a
single culprit; Theorem~\ref{thm:byzantine} bounds worst-case \emph{colluding-coalition}
identification (block-spark $>2f$), a stronger adversary. (c)~They treat clock \emph{jumps} but
do not characterize the common offset and rate as an \emph{unobservable gauge with a holdover
floor} (Corollary~\ref{cor:holdover}), nor the coupled datum$\oplus$timescale unification
(Proposition~\ref{prop:unification}), nor the provider common/differential split. The deep
connection we concede: EDM/rigidity is intrinsically invariant to rigid transforms, so those
methods \emph{exploit} the spatial datum gauge as a tool (hence needing no precise ephemeris),
while we \emph{characterize} that invariance as the fault-hiding subspace and extend it to the
timescale. Our results and theirs are thus complementary.

\emph{Versus secure state estimation \citep{fawzi2014,shoukry2016,candes2005,donoho2003}.} The
sparse-observability / $2f$ bound is entirely theirs. Theorem~\ref{thm:byzantine} is its
geometric instantiation---the block-spark of the peer column-blocks of a physical lunar
$\Gmat$ carrying the coupled datum gauge. We claim the \emph{instantiation and its lunar
consequences} (the deficiency-of-five self-blindness, the provider split), not the bound.

\emph{Versus LiAISON \citep{hillborn2007}.} LiAISON obtains absolute state from crosslinks by
exploiting orbital dynamics over an arc; our analysis is the kinematic snapshot, whose gauge
is eight (Section~\ref{sec:scope}). The dynamical regime is the complementary way to reduce
that gauge, and we do not claim a dynamics-free absolute frame.

\emph{What is ours.} The coupled frame$\oplus$timescale gauge instantiated in fault
observability (the undetectable subspace is the measurement-space \emph{face} of the same rank
deficiency that produces the datum$\oplus$timescale gauge in parameter space); the geometric
secure-\emph{identification} bound for colluding peers;
the P1/P3/P4 gauge-unification; and the provider-fault taxonomy. Not ours: parity RAIM, the
MDB, the PL, the spark bound, the Byzantine model, and rigidity-based detection---all cited.

\subsection{Relation to the companion arc}
\label{sec:companions}

This is the fifth and closing paper of a program built on one open engine. P1 classifies the
spatial datum defect \citep{baweja2026datum}; P2 the cross-provider interoperability floor
\citep{baweja2026interop}; P3 the common-mode class an integrity monitor cannot see
\citep{baweja2026commonmode}; P4 the coupled frame$\oplus$timescale gauge
\citep{baweja2026gauge}. The present paper shows that these are not four separate gauges but
the fault-observability faces of one rank deficiency: the spatial datum (P1) and the temporal
gauge (P4) are the parameter-space $\Nsp(\Gmat)$; the common-mode class (P3) is a subset of
the measurement-space $\rng(\Gmat)$; and the interoperability floor (P2) is the common,
externally-tied part of the provider channel. Proposition~\ref{prop:unification} is the
statement that closes the loop. The five share the same firewall---exactly one externally
validated row per paper---and the same open substrate.

\subsection{Limitations}
\label{sec:limitations}

The firewall of Section~\ref{sec:anchor} is also the statement of limitations; we make them
explicit.
\begin{itemize}
  \item \textbf{Kinematic snapshot scope.} $\Gmat$ is the instantaneous geometry Jacobian;
  the gauge of eight is the snapshot datum defect. Orbital dynamics over an arc
  \citep{hillborn2007} can partially collapse $\Nsp(\Gmat)$ and is not modelled here. Our
  results are the geometry regime.
  \item \textbf{Representative geometry.} The eight-node network is a construction, not a
  flown constellation. The block-spark tolerances, the slope, the parity crossing, and the
  provider counts depend on it; the proven \emph{formulas} and structural results do not.
  \item \textbf{Counts are validated as computations, not as guarantees.} The Byzantine $f$,
  the provider-count row, and the protection crossing are outputs of the validated pipeline;
  read as operational flight tolerances they are \Modelled{}. There is no flight data.
  \item \textbf{First-order linearization.} The state, rows, and normal matrix are first
  order about a nominal; second-order geometric effects and non-Gaussian errors are outside
  the model.
  \item \textbf{No flight, certification, or TRL claim.} The alert limit is hypothetical, the
  holdover magnitude is a modelled oscillator drift, and nothing here is certified or endorsed
  by any agency.
\end{itemize}
None of these caveats touches the \Proven{} content---the projector survival, the
undetectable characterization, the unification, the block-spark bound, the holdover floor, and
the protection gap---or the single \Validated{} pipeline reproduction.

\subsection{Venue and standards relevance}
\label{sec:venue}

The results speak to autonomous integrity for lunar PNT (a natural fit for \emph{NAVIGATION}
or IEEE TAES, and for ION GNSS+ 2026) and to the current lunar reference-system and lunar-time
standardization---both the IAU~2024 coordinate time (Lunar Coordinate Time, TCL) and the
realized Coordinated Lunar Time (LTC) motivated by the 2024 US OSTP memorandum
\citep{iau2024ltc,fienga2024moonlight,lnis2023}. The
design laws are stated as relevance to those efforts---an autonomous architecture needs
engineered redundancy, distributed measurement coverage, an external tie for the common
classes, and an acknowledgement of the protection gap---not as certified requirements or as
any form of agency endorsement.
